\section{Introduction}

A financial investigator queries a cloud-hosted transaction graph: find all 
accounts reachable within three hops from a set of flagged entities. The data 
is encrypted, and the query executes inside a Trusted Execution Environment 
(TEE) --- hardware that prevents the cloud provider from reading any data 
values. Yet an adversary watching memory access patterns during query execution 
can still reconstruct which accounts were touched, infer node degrees, and 
trace the investigator's interest, all without decrypting any data. 

TEEs provide hardware-level isolation that protects data \emph{values} from the 
cloud provider, and have emerged as the leading approach for confidential cloud 
computation. But they do not protect \emph{access patterns}. An adversary 
observing which memory addresses a query touches, in what order and with what 
frequency, can infer the graph's structure, reconstruct query predicates, and in 
some cases recover significant portions of the underlying 
data~\cite{islam2012access, grubbs2019learning, kellaris2016generic, 
kornaropoulos2019data, cash2015leakage, demertzis2020seal, grubbs2018pump, 
gui2019encrypted, lacharite2018improved, oya2021hiding, poddar2020practical, 
zhang2016all, oya2022ihop}. For regulated industries handling financial or 
healthcare data, this residual leakage may lead to compliance failure.

The solution is \emph{oblivious} computation: algorithms whose memory access 
patterns and control flow are provably independent of the input data. The 
challenge is executing this efficiently for property graph queries. Graph 
traversal is structurally leaky by nature: multi-hop queries, which follow 
chains of node-edge-node patterns across the graph, expose intermediate fanout 
at every step. A standard hash join reveals node degrees through its probe 
pattern. While Oblivious RAM 
(ORAM)~\cite{goldreich1987towards, stefanov2013path} can be used to hide 
individual accesses, it introduces logarithmic overhead per operation, making 
it impractical for multi-hop query workloads at scale. Prior approaches ~\cite{chamani2023graphos, 
feng2025enabling} address oblivious graph \emph{algorithms}, such as BFS, DFS, 
minimum spanning tree, but these differ fundamentally from the selective, 
predicate-driven queries that property graph databases execute in practice. 
Adapting oblivious relational join 
operators~\cite{krastnikov2020efficient, zheng2017opaque, 
eskandarian2019oblidb, arasu2013oblivious, mavrogiannakis2025obliviator, 
reichert2024menhir} to graph queries is possible, but a naive decomposition 
into binary joins leaks intermediate output sizes --- themselves sensitive 
signals about the graph structure being queried.


The closest prior works on fully oblivious multi-way joins operate at the
purely relational level. Wei and Kerschbaum~\cite{ruidi2025oblivious} (WK) 
compute oblivious multi-way joins while hiding intermediate result sizes, 
providing strong privacy guarantees for acyclic join queries. Hu and 
Wu~\cite{hu2025optimal} achieve worst-case optimal oblivious joins matching the 
AGM bound up to a logarithmic factor, without relying on ORAM or cryptographic 
assumptions. Yet, both schemes are designed for general relational engines,
unaware of the structural properties specific to 
property graph workloads. This leaves substantial performance untapped.

Moreover, these schemes --- along with other oblivious join-then-filter 
approaches such as Obliviator~\cite{mavrogiannakis2025obliviator} --- treat 
filtering as a post-processing step applied after the join materializes its 
full output. This is both expensive, since the unfiltered join result can be 
orders of magnitude larger than the filtered one, and privacy-weakening, since 
the size of the unfiltered intermediate is revealed to the adversary.

\emph{Our insight.} Property graph databases exhibit two structural properties 
that we exploit. First, edges are stored in 
both \emph{forward} and \emph{reverse} directions, sorted by source and 
destination node ids respectively, a convention common in production graph 
systems~\cite{janusgraph, feng2023kuzu, deutsch2019tigergraph, 
sun2015sqlgraph}. A one-hop traversal $n_i \rightarrow e \rightarrow n_j$ can 
therefore be evaluated from \emph{both ends simultaneously} and the partial 
results merged, enabling parallelism. Second, graph 
databases support stored-procedure-like query 
templates~\cite{deutsch2019tigergraph, neo4j, janusgraph, bebee2018amazon}: 
the schema of nodes and edges involved in a query is fixed ahead of time, with 
only predicate values supplied at query time. This enables an offline 
preprocessing phase to prepare node and edge tables, so that online 
execution issues fewer oblivious operations.
Additionally, graph queries are predicate-driven, selecting nodes 
and edges that satisfy specific property constraints. Hence, filtering should 
be an integral part of every processing step, and not a separate phase. 


These observations motivate a new primitive: an efficient \emph{one-hop 
operator} that processes $n_i \rightarrow e \rightarrow n_j$ traversals 
obliviously, exploiting bidirectional edge layout and the offline/online 
processing. Most multi-hop queries decompose into a set of one-hop 
subqueries, which are \emph{independent} and can execute in parallel. The 
residual join across decomposed parts is handled by the WK~\cite{ruidi2025oblivious} multi-way join, 
which preserves its privacy guarantees on intermediate sizes throughout.
We 
build on WK rather than Hu and Wu's scheme~\cite{hu2025optimal} because WK provides a practical 
implementation amenable to systems integration, whereas \cite{hu2025optimal}'s 
contribution is primarily theoretical, whose 
constant factors and cache assumptions have not yet been validated in a 
practical system.

\emph{Contributions.} We present \system, an oblivious property graph database 
built around three contributions:

\begin{itemize}
    \item \textbf{One-hop oblivious operator.} A new operator that exploits 
    bidirectional edge storage and offline/online decomposition to process 
    single-hop traversals significantly more efficiently than a general 
    oblivious join.

    \item \textbf{In-place oblivious filtering.} A mechanism that fuses 
    predicate filtering directly into the join computation, avoiding 
    materializing the unfiltered intermediate result. Unlike existing 
    join-then-filter schemes~\cite{ruidi2025oblivious, 
    mavrogiannakis2025obliviator}, \system never reveals any size except the 
    final filtered output, improving both performance and privacy.
    
    \item \textbf{Query decomposition framework.} A rewrite strategy that 
    decomposes multi-hop graph queries to maximize parallelism across one-hop 
    operators, using WK as a black-box combiner for the residual join, while 
    preserving identical privacy guarantees throughout.

    \item \textbf{Empirical evaluation.} An evaluation on a synthetic banking 
    transaction dataset demonstrating significant end-to-end speedups over 
    non-decomposed oblivious execution for representative 4-hop chain and 
    branch query patterns.
\end{itemize}

Together, these contributions demonstrate that treating property graph 
databases as first-class citizens --- rather than as relational engines 
operating on graph-shaped data --- yields meaningful performance gains in the 
oblivious setting without weakening the underlying security guarantees.